% =====================================================================
% I. Introduction  (틀 — 내용 새로 작성)
% =====================================================================
\section{Introduction}\label{sec:intro}

% [작성 가이드]
%   - 첫 문단은 \IEEEPARstart 드롭캡으로 시작(IEEE 관례). 첫 글자 + 첫 단어.
%   - 배경/동기 → 문제 정의 → 기존 접근의 한계 → 본 논문의 접근 → 기여 목록
%     → 논문 구성(organization) 순.

\IEEEPARstart{X}{xxx} \todo{첫 문단: 연구 배경과 동기.}

\todo{문제 정의와 기존 시스템이 깨지는 지점(가정/한계).}

% --- 기여 목록 (원하는 개수만큼 \item) --------------------------------
\noindent The contributions of this paper are as follows:
\begin{itemize}
  \item \todo{기여 1}
  \item \todo{기여 2}
  \item \todo{기여 3}
\end{itemize}

% --- 논문 구성 --------------------------------------------------------
\todo{The rest of this paper is organized as follows. Section~\ref{sec:related} ...
Section~\ref{sec:disc} concludes.}
